% original p. 342
\sgasection{7}{APPENDIX: Existence of regular elements over finite fields}
\begin{center}by J.-P. Serre\end{center}
\setcounter{footnote}{0}
Throughout what follows, $k$ denotes a finite field and $\bar k$ its
algebraic closure; the Galois group of $\bar k/k$ is denoted by
$\mathscr G$; recall that if $q=\operatorname{Card}(k)$,
$\mathscr G$ is topologically generated by the ``Frobenius'' element
$F:x\mto x^q$.

One proposes to prove the following theorem\footnote{This theorem
is due to Chevalley. The writer wishes to express his gratitude
to American Express which, by mislaying a trunk of Chevalley's
manuscripts, obliged him to reconstruct the proof.}:
\begin{theorem}{}\label{XIV.Appendix.Theorem}
Let $G$ be an adjoint semisimple group defined over $k$, and
$\mathfrak g$ its Lie algebra. The $k$-Lie algebra
$\mathfrak g(k)$ contains a regular element.
\end{theorem}
\begin{enoncerm}{Remarks}{}
(1) It is well to recall that ``adjoint'' means that the center
of $G$ is trivial (as a group subscheme of $G$). In view of the
Chevalley seminar, this also means that if $T$ is a maximal
torus of $G$, the group $X(T)$ of characters of $T$ (defined
over $\bar k$) is generated by the set $R$ of roots. It follows
in particular that the rank of the Lie algebra $\mathfrak g$
is equal to the dimension of $T$ (i.e. to the rank of $G$).

(2) The writer does not know whether the hypothesis that
$G$ is adjoint is indispensable.
\end{enoncerm}

\begin{lemma}{1}\label{XIV.1}
It suffices to prove the theorem when $G$ is geometrically simple.
\end{lemma}
(One says that $G$ is geometrically simple if $G\otimes\bar k$
contains no normal group subscheme smooth over $\bar k$ apart
from $G$ and $\{e\}$; an equivalent condition: the associated
root system $R$ is irreducible.)

One may first suppose $G$ indecomposable over $k$. The group
$G\otimes\bar k$ is then a product of geometrically simple
components $S$ which are permuted transitively by the Galois
group $\mathscr G$. If $\mathscr H$ is the subgroup of
$\mathscr G$ fixing one of these components $S$, this component
is defined over the field $K$ corresponding to $\mathscr H$
(i.e. comes by extension of scalars from a subscheme of
$G\otimes K$), and a standard argument shows that
$G=\prod_{K/k}(S)$ (i.e. $R_{K/k}(S)$, for readers accustomed
to Weil's notation). Likewise, the Lie algebra $\mathfrak g$
of $G$ is identified with $\prod_{K/k}\mathfrak s$, where
$\mathfrak s$ is the Lie algebra of $S$. If the theorem is
true for $S$, there exists $x\in\mathfrak s(K)=\mathfrak g(k)$
which is regular in $\mathfrak s$; one then easily verifies
that it is regular in $\mathfrak g$. \hfill Q.E.D.

From now on, one supposes $G$ geometrically simple and chooses
a maximal torus $T$ of $G$ (one knows that this is possible).
Denote by $X$ the character group of $T\otimes\bar k$,
by $R$ its root system, by $W$ its Weyl group, and by $E$
the group of automorphisms of $X$ preserving $R$. The group
$W$ is a normal subgroup of $E$.

If $T'$ is another torus of $G$, denote by $X',R',W',E'$
the corresponding character group, root system,\ldots.
If one chooses $y\in G(\bar k)$ such that
$y\cdot(T\otimes\bar k)\cdot y^{-1}=T'\otimes\bar k$,
one may identify $X',R',W',E'$ with $X,R,W,E$ by means
of $\operatorname{int}(y)$. Changing $y$ changes this
identification by an automorphism of $X$ corresponding
to an element of $W$.

The canonical generator $x\mto x^q$ of $\mathscr G$
acts on $X'$ leaving $R'$ stable; it therefore defines
an element $f_{T'}$ of $E'$. One will in particular
set $f=f_T$. When one identifies $E'$ with $E$ as
just described, the element $f_{T'}$ becomes an
element $f'$ of $E$, defined up to replacement by
$wf'w^{-1}$, with $w\in W$. One will compare this
element to $f$:

\begin{lemma}{2}\label{XIV.2}
One has $f'\equiv f\pmod W$; conversely, every
$f'\in E$ satisfying this condition can be
obtained from a maximal torus $T'$ of $G$.
\end{lemma}
Set $\bar T=T\otimes\bar k$, $\bar T'=T'\otimes\bar k$,
and let $y\in G(\bar k)$ be such that
$y\bar Ty^{-1}=\bar T'$; since
$y^q\bar Ty^{-q}=\bar T'$, one concludes that
$n=y^{-1}y^q$ belongs to the normalizer
$N(\bar T)$ of $\bar T$. The effect of $f'$
on the points of $\bar T(\bar k)$ is then
as follows:
\[
 f'(t)=y^{-1}(yty^{-1})^qy
 =y^{-1}y^qt^qy^{-q}y=nt^qn^{-1}.
\]
If $w\in W$ is the element defined by
$n$, this shows that $f'$ and $f\circ w$
have the same effect on the points of
$\bar T$, hence also on its characters,
and one has $f'=f\circ w$, whence
$f'\equiv f\pmod W$. Conversely, if
$w\in W$ is given, one represents it
by an element $n\in N(\bar T)(\bar k)$;
by a classical theorem of Lang, one
may write $n$ in the form $n=y^{-1}y^q$,
with $y\in G(\bar k)$; the torus
$\bar T'=y\bar Ty^{-1}$ is then
defined over $k$, and the preceding
calculation shows that the corresponding
$f'$ is equal to $f\circ w$.

\begin{lemma}{3}\label{XIV.3}
Let $X,R,W,E$ be as above ($R$
being irreducible), and let
$\varphi\in E/W$. There then exists
an element $f\in E$ representing
$\varphi$, and a family
$\theta_1,\ldots,\theta_n$ of roots
having the following two properties:
\begin{enumerate}[label=\textup{(\arabic*)}]
\item $(\theta_1,\ldots,\theta_n)$ is a basis of $X$.
\item $R$ is the union of the orbits of the $\theta_i$
under the powers of $f$ (i.e. every $a\in R$ is
written $a=f^m\theta_i$, with suitable $m$ and $i$).
\end{enumerate}
\end{lemma}
The proof will be given a little later.

Here is now a linear algebra lemma:
\begin{lemma}{4}\label{XIV.4}
Let $V$ be a vector space over $k$, and
$(\theta_1,\ldots,\theta_n)$ a basis of
the dual of $V\otimes\bar k$. There
exists $x\in V$ such that $\theta_i(x)\ne0$
for every $i$.
\end{lemma}
Let $V^*$ be the dual of $V$, and
$\bar W_i$ the subspace of
$V^*\otimes\bar k$ generated by
the conjugates of $\theta_i$;
this subspace is defined over
$k$, i.e. of the form
$W_i\otimes\bar k$. The obvious
map $W_1\otimes\cdots\otimes W_n\to\bigwedge^nV^*$
is nonzero (otherwise its
extension to $\bar k$ would
also be zero, which is absurd
since $\theta_i\in\bar W_i$,
and $\theta_1\wedge\cdots\wedge\theta_n\ne0$).
There therefore exists a
basis of $V^*$ consisting
of elements $u_i\in W_i$.
Let $x\in V$ be such that
$u_i(x)\ne0$ for every $i$
(for example $u_i(x)=1$).
The element $x$ answers
the question, for if one
had $\theta_i(x)=0$, the
conjugates of $\theta_i$
would also vanish at $x$,
and the same would hold
for $u_i$, which is not
the case.

\par\medskip\noindent\emph{End of the proof of the theorem.}
Combining lemmas 2 and 3,
one may choose a torus $T$
whose element $f$ satisfies
the properties of lemma 3.
If $Y$ is the dual of $X$,
the Lie algebra
$\mathfrak t(\bar k)$ of
$\bar T$\setcounter{footnote}{8}\nde{$T$
has been changed to $\bar T$.}\setcounter{footnote}{1}
is canonically identified
with $Y\otimes\bar k$,
and this operation is
compatible with the action
of the Galois group (the
latter acting on
$Y\otimes\bar k$ by its
action on $Y$ and on
$\bar k$). Let
$V=\mathfrak t(k)$ be
the Lie algebra of $T$
over $k$. An element
$x\in V$ is regular if
and only if it is
annihilated by no root
$\alpha\in R$, or
rather by none of the
linear forms
$\bar\alpha\in V^*\otimes\bar k$
defined canonically by
the $\alpha\in R$.
By lemma 4, one may
find such an $x$
which is annihilated
by none of the roots
$\theta_i$; but every
root is conjugate to
a $\theta_i$ (this is
what condition (2)
of lemma 3 expresses);
it follows that $x$
is annihilated by no
root, and it is
indeed a regular
 element.
% typo: de Lie de de T -> de Lie de T.

\par\medskip\noindent\emph{Proof of lemma 3.}
It relies on the
properties of
\emph{Coxeter transformations}.
Let us briefly recall
what this consists of:

Let $a_1,\ldots,a_n$
be a simple system
of roots of $R$,
and for every $i$
let $r_i$ be the
reflection corresponding
to $a_i$. Set:
\[
 c=r_1\cdots r_n.
\]
One has $c\in W$;
of course, the element
$c$ depends on the
choice of the simple
system $(a_1,\ldots,a_n)$
as well as on the
order of the $a_i$;
however, one proves
that its conjugacy
class depends on
none of these choices.
One calls $c$ the
\emph{Coxeter element}
of the system under
consideration. One
proves (we shall
admit it) that $c$
does not have $1$
as eigenvalue.

\begin{lemma}{5}\label{XIV.5}
Set $\theta_i=r_n\cdots r_{i+1}(a_i)$. Then:
\begin{enumerate}[label=\textup{(\alph*)}]
\item The $\theta_i$ form a basis of the group $X$
generated by the roots.
\item One has $\theta_i>0$ and $c(\theta_i)<0$ for
every $i$.\footnote{For the order relation defined
by $(a_1,\ldots,a_n)$.}
\item Every root $a\in R$ such that $a>0$ and
$c(a)<0$ is equal to one of the $\theta_i$.
\item $R$ is the union of the orbits of the
$\theta_i$ under the powers of $c$.
\end{enumerate}
\end{lemma}
It is clear that $\theta_i$ is of the form
$a_i+\sum_{j>i}m_{ij}a_j$, with $m_{ij}\in\ZZ$,
which proves (a). Assertions (b) and (c)
are consequences of the following observation:
the reflection $r_i$ preserves the sign of
every root distinct from $\pm a_i$, and
changes the sign of $\pm a_i$.

Finally, for (d), one observes that an
orbit of $c$ cannot consist entirely
of positive (resp. negative) roots,
for by taking the sum of these roots
one would find a nonzero element of
$X$ invariant under $c$, and one has
been willing to admit that $c$ does
not have $1$ as eigenvalue. There
is therefore necessarily in every
orbit an element $a>0$ such that
$c(a)<0$, and one applies (c).

\begin{remark*}
The preceding proof has been sketched
only to facilitate the reader's task;
one could have confined oneself to
referring to the canonical texts on
Coxeter (cf., for example, Koszul,
S\'eminaire Bourbaki, 1959/1960,
expos\'e 191). These texts contain
other results: the orbits of $c$
all have the same number $h$ of
elements, and each contains only
one $\theta_i$. In particular,
$h=\operatorname{Card}(R)/n$.
\end{remark*}

Let us now return to the proof
of lemma 3. Distinguish three cases:

\noindent (1) \emph{The given element
$\varphi\in E/W$ is equal to the
identity element.}

One must then take $f\in W$;
one chooses $f=c$; this works
by lemma 5.

\noindent (2) \emph{The system $R$
has type $A_n$, with even $n\geq2$,
and $\varphi$ is the unique
nontrivial element of $E/W$.}

One knows that $-1\notin W$;
one then takes $f=-c$. A simple
calculation shows that $c$ has
order $h=n+1$; its order is
therefore odd. If $a\in R$
is any root, one has
$a=c^m\theta_i$ for a pair
$(m,i)$; cf. lemma 5; adding
$h$ to $m$ if necessary,
one may suppose $m$ even,
and one sees that one then
has $a=(-c)^m\theta_i=f^m\theta_i$.
The orbits of the $\theta_i$
therefore indeed fill $R$.

\noindent (3) \emph{The element
$\varphi\in E/W$ is nontrivial,
and $R$ has one of the
following types:}
\begin{center}
\begin{tabular}{@{}ll@{}}
\rule{0pt}{3em}$A_n,\ n\text{ odd}\geq3$ &
$\xymatrix@C=1.1em{\bullet\ar@{-}[r]\ar@{<->}@/^1.6pc/[rrrr]&
\bullet\ar@{-}[r]\ar@{<->}@/^1pc/[rr]&\bullet\ar@{-}[r]&\bullet\ar@{-}[r]&\bullet}$\\[3em]
\rule{0pt}{3em}$D_4$ &
$\xymatrix@C=1.2em@R=1em{
 &\bullet\ar@{->}@/_1.4pc/[dd]&&\\
 &&\bullet\ar@{-}[ul]\ar@{-}[dl]\ar@{-}[r]&\bullet\ar@{->}@/_1.4pc/[ull]\\
 &\bullet\ar@{->}@/_1.4pc/[urr]&&}$\\[3em]
\rule{0pt}{3em}$D_n,\ n\geq5$ &
$\xymatrix@C=1.1em@R=.7em{
 &&&&\bullet\ar@{<->}@/^1pc/[dd]\\
 \bullet\ar@{-}[r]&\bullet\ar@{-}[r]&\bullet\ar@{-}[r]&\bullet\ar@{-}[ur]\ar@{-}[dr]&\\
 &&&&\bullet}$\\[3em]
\rule{0pt}{3em}$E_6.$ &
$\xymatrix@C=1.1em@R=1.1em{
 \bullet\ar@{-}[r]\ar@{<->}@/^1.6pc/[rrrr]&
 \bullet\ar@{-}[r]\ar@{<->}@/^1pc/[rr]&
 \bullet\ar@{-}[r]\ar@{-}[d]&\bullet\ar@{-}[r]&\bullet\\
 &&\bullet&&}$
\end{tabular}
\end{center}
(A glance at the classification
shows that these are indeed
all the cases (with $A_n$,
even $n$) where $E/W$ is
nontrivial.)

Let $S$ be a simple system
of roots, which we do not
number for the moment.
One knows that $E$ is the
semidirect product of $W$
and the group $\Psi$ of
permutations of $S$ which
leave the Cartan matrix
invariant (or the Dynkin
diagram; it is the same
thing). The group $E/W$
is thus identified with
$\Psi$, and in particular
$\varphi$ corresponds to
an element $\psi\in\Psi$.
One observes by inspecting
the Dynkin diagrams (cf.
the figures above) that
every orbit of $\psi$ in
$S$ consists of pairwise
orthogonal roots (i.e.
not joined in the diagram).
[Note that this would
not be the case for
$A_n$ (even $n$), which
obliged us to treat this
case separately.] If
$\sigma$ is such an
orbit, the reflections
$r_a$, $a\in\sigma$,
commute with one another;
their product will be
denoted by $\rho_\sigma$.
It is clear that
$\rho_\sigma$ commutes
with $\psi$.

This being so, let us
choose a total order
on $S$ such that every
orbit is a segment
for this order relation;
this amounts to numbering
the elements of $S$:
$S=\{a_1,\ldots,a_n\}$.
Let $\theta_1,\ldots,\theta_n$
be the roots defined
above, and
$c=r_1\cdots r_n$
the corresponding
Coxeter element.
The element $c$ is
a product of the
$\rho_\sigma$, the
$\sigma$ being arranged
in a certain order;
it follows that it
commutes with $\psi$.
One then sets $f=c\psi$.
One observes in
addition that $\psi$
permutes the $\theta_i$
among themselves.
Indeed, one has
$\psi(\theta_i)>0$
since $\theta_i>0$,
and
$c(\psi(\theta_i))=\psi(c(\theta_{i'}))<0$,
hence (lemma 5, (c)),
$\psi(\theta_i)$ is
equal to a $\theta_j$.
% typo?: kept as printed: theta_{i'} inside psi(c(...)).
It is now immediate
that $f=c\psi$
answers the question.
Indeed, if $a\in R$,
one has $a=c^m\theta_i$
for a pair $(m,i)$,
whence
$a=f^m(\psi^{-m}\theta_i)=f^m\theta_j$
for a certain $j$.
\hfill Q.E.D.

\begin{remark*}
One can prove that,
except in case (2),
every orbit of $f$
has exactly
$h=\operatorname{Card}(R)/n$
elements and contains
one and only one
$\theta_i$. In case
(2), some of the
$\theta_i$ are
unnecessary.
\end{remark*}
% Source PDF p. 36 is blank.
